%=============================================================================
% MODULE 08: CONCLUSION
%=============================================================================
% Frontiers in Public Health — Concise conclusion, no new data
%=============================================================================

\section{Conclusions}

This study demonstrates significant algorithmic bias in diabetes prediction 
models across demographic groups. The SVM model achieved the best 
discriminative performance (F1-Score: 0.8313, AUC-ROC: 0.8794) but still 
exhibited notable fairness concerns, particularly at the intersection of 
race and gender. All four evaluated models showed higher false positive rates 
for minority patients, indicating systematic over-diagnosis bias.

The intersectional analysis revealed that minority females experience 
amplified disadvantage, with a true positive rate of only 0.714 compared 
to 0.833 for majority groups. This finding underscores the limitations of 
single-attribute fairness analysis and highlights the importance of examining 
bias at the intersection of multiple marginalized identities.

Our results underscore the urgent need for:

\begin{enumerate}[leftmargin=*]
    \item Systematic bias assessment in clinical ML models prior to deployment, 
          using multiple fairness metrics and intersectional analysis.
    \item Fairness-aware model development that explicitly incorporates equity 
          objectives alongside performance optimization.
    \item Continuous monitoring for bias in deployed systems, with automated 
          alerts when disparities exceed acceptable thresholds.
    \item Transparent reporting of performance across demographic groups, 
          enabling clinicians and patients to make informed decisions about 
          AI-assisted care.
\end{enumerate}

As artificial intelligence becomes increasingly integrated into healthcare 
delivery, ensuring algorithmic fairness is not merely a technical challenge 
but a moral imperative. Only through deliberate, systematic attention to 
equity can we realize the promise of AI for all patients, regardless of 
their demographic background. The path forward requires collaboration 
across disciplines---including computer science, clinical medicine, 
ethics, and public health---to develop and deploy AI systems that 
advance rather than undermine health equity.
